% =====================================================================
\lezione{02/11/2026}{Anomaly detection: concetti, one-class SVM e k-means}
  {Prof.~Danilo Giordano}
  {anomaly, outlier e novelty detection; perché in cybersecurity; complessità e tipi di anomalie (point,
   contextual, collective); sfide per AI/ML; supervised vs unsupervised; score e label; shallow learning;
   SVM e kernel trick; one-class SVM e parametro $\nu$; clustering e k-means; SSE e metodo del gomito}
\label{lez:2026-11-02}
% =====================================================================

\begin{nota}[Fonte]
Trascrizione automatica della registrazione, revisionata, confrontata con le slide
\emph{5-Supervised-vs-Unsupervised-shallow} (pp.~1--33).
\end{nota}

\begin{verifica}
Tra la lezione del 20/10 e questa ci sono state attività non registrate (le due sessioni del Lab~2 e
le eventuali lezioni di quei giorni). Il diario non ne riporta il contenuto.
\end{verifica}

\section{Sintesi}
Prima lezione dell'ultimo blocco del Prof.~Giordano, dedicato all'\textbf{anomaly detection}; dopo
inizierà l'NLP con il Prof.~Boffa. Il punto di partenza: i classificatori falliscono quando le label
mancano, sono molto sbilanciate o inaffidabili, o quando si vogliono scoprire \textbf{novità}. La lezione
definisce il problema, le sue complessità e le sfide per il ML, confronta l'approccio supervisionato e
quello non supervisionato e inizia gli algoritmi \emph{shallow}: SVM, \textbf{one-class SVM} e
\textbf{k-means}. Si ferma alla scelta di $K$; l'uso di k-means per l'anomaly detection è rimandato alla
lezione successiva. Approfondimento nel capitolo~\ref{cap:anomaly}.

\section{Avvisi pratici}
\begin{itemize}
  \item Il docente chiede un riscontro sul Lab~2 (``difficile?'') e lo difende come utile, anche se
    impegnativo.
  \item La parte su k-means per l'anomaly detection riprende ``domani o mercoledì'': quella lezione non è
    nelle registrazioni (vedi \rif{2026-11-09}{09/11}).
\end{itemize}

\section{Contenuti}

\subsection{Che cos'è e perché serve}
L'anomaly detection cerca \textbf{pattern che deviano dal comportamento atteso}. Va distinta la
\textbf{novelty detection}: un pattern nuovo può essere la \textbf{nuova normalità} (l'uso della VPN
durante il COVID). In cybersecurity le anomalie segnalano attività malevole (zero-day, intrusioni,
esfiltrazione dopo un movimento laterale), guasti e configurazioni errate, frodi bancarie. Serve perché
per uno \textbf{zero-day} non può esistere un classificatore, perché gli attaccanti si adattano alle
regole (gatto e topo) e per avere un \textbf{early warning} prima del danno.

\subsection{Complessità, tipi di anomalie e sfide}
Le complessità sono quattro: \textbf{unknowingness} (le anomalie sono ignote finché non accadono),
\textbf{rarità e sbilanciamento} (e sbagliare su un'anomalia costa di più), \textbf{classi eterogenee}
(lo stesso tipo di attacco si presenta in forme diverse, es.\ DNS tunnelling e HTTP/2 push) e
\textbf{pattern che evolvono}. I tipi di anomalia sono tre:
\begin{itemize}
  \item \textbf{point}: un'istanza anomala da sola (login da un paese mai visto);
  \item \textbf{contextual}: anomala solo dato il contesto (un backup di due ore alle 3 di domenica);
  \item \textbf{collective}: un insieme di istanze normali che insieme formano l'anomalia (tentativi di
    login ognuno sotto la soglia, ma in sequenza).
\end{itemize}
Le sfide per AI/ML: recall bassa e molti falsi positivi (se la maggior parte degli allarmi è falsa,
nessuno si fida più del modello), dati ad alta dimensionalità e dipendenti, poche label, rumore, anomalie
complesse e bisogno di \textbf{spiegazioni}. Approccio \textbf{supervisionato}: accurato con buone label
ma non generalizza a nuove anomalie; \textbf{non supervisionato}: assume che il normale domini, trova
minacce nuove ma ha più falsi positivi ed è meno interpretabile. L'output di un algoritmo può essere uno
\textbf{score}, che permette di ordinare le istanze e di lavorare a budget (esempio Intesa Sanpaolo:
le prime 10.000 del mese), oppure una \textbf{label}, il cui numero si controlla con gli iperparametri.

\begin{esame}[un'anomalia non è per forza un attacco]
Il docente torna più volte sul contesto: un rilevatore che blocca subito tutte le VPN perché il loro uso
è ``anomalo'' fermerebbe tutti i lavoratori da remoto. Saper distinguere anomalia, novità e attacco, e
riconoscere point, contextual e collective anomaly negli esempi.
\end{esame}

\subsection{Shallow learning: SVM e one-class SVM}
Gli algoritmi shallow sono interpretabili, vanno bene sui dati tabellari e sulle anomalie puntuali e
danno una \textbf{baseline}; sono deboli sulle anomalie contestuali e collettive, che richiedono feature
costruite a mano (es.\ tentativi falliti aggregati su mezz'ora). La \textbf{SVM} (1995) cerca l'iperpiano
che separa le classi \textbf{massimizzando il margine}; il \textbf{kernel trick} porta i dati in uno spazio
in cui diventano separabili (lineare, polinomiale, \textbf{RBF}, il più usato). La \textbf{one-class SVM}
non riceve label: costruisce un confine attorno ai dati normali, e ciò che sta fuori è anomalo. Il
parametro $\bm{\nu}$ è la \textbf{frazione attesa di outlier}. Si può usare in due modi: addestrata
\textbf{solo su dati normali} (``supervisionata'' tra virgolette: le label servono solo a scegliere i
dati, $\nu$ piccolo) o su \textbf{tutti i dati} senza label (serve una stima della frazione di anomalie).

\subsection{Clustering e k-means}
Nel clustering le anomalie sono i \textbf{cluster piccoli e lontani}, i punti \textbf{lontani dal
centroide} del proprio cluster o i punti \textbf{senza cluster}. K-means è partizionale e basato su
centroidi: ogni punto appartiene a un cluster, funziona con cluster sferici di dimensioni simili.
Dipende dall'\textbf{inizializzazione} casuale, quindi si esegue più volte con lo stesso $K$ e si tiene il
risultato migliore. $K$ si sceglie con l'\textbf{SSE} e il \textbf{metodo del gomito}, restando nella
zona del gomito: più cluster abbassano l'SSE ma vanno studiati uno per uno, senza label.

\begin{esame}[SSE che risale]
Il docente chiede di ricordarlo: se nel grafico SSE contro $K$ compare un picco verso l'alto,
probabilmente un'inizializzazione sfortunata ha lasciato un \textbf{cluster vuoto}. Si riesegue k-means.
\end{esame}

\section{Domande aperte e punti da verificare}
\begin{verifica}
\begin{itemize}
  \item La slide sulla one-class SVM dice che con $\nu$ più piccolo si hanno meno outlier e confini
    ``più stretti'': ``meno outlier'' è corretto, ma lasciando fuori meno punti la regione normale si
    allarga. Da chiarire con il docente.
  \item Contenuto delle lezioni e dei lab non registrati tra il 20/10 e il 02/11: non ricostruito.
\end{itemize}
\end{verifica}
